% CONCLUSÕES

\chapter{CONCLUSÕES}
\label{chap:conclusoes}

Este Trabalho de Conclusão de Curso investigou o impacto do uso do ChatGPT na compreensão de estudantes da disciplina de Algoritmos e Estruturas de Dados, com foco nos conteúdos de Heap e Tabelas Hash. O estudo foi conduzido em contexto de graduação em Engenharia da Computação, adotando um desenho experimental com grupos cruzados e questionários objetivos aplicados ao final de cada atividade prática.

Em relação ao objetivo geral proposto, os resultados indicam que o uso do ChatGPT durante as atividades práticas não produziu diferença estatisticamente significativa na compreensão dos conteúdos, medida pelos questionários de Heap e Tabelas Hash. Na comparação agregada entre as condições com e sem uso da ferramenta, a diferença entre as médias foi de apenas 0,15 ponto, o teste de Mann-Whitney U não permitiu rejeitar a hipótese $H_{0}$ ($p = 1{,}000$) e o Cliff's Delta indicou efeito negligível ($\delta = -0{,}003$). Análises separadas por tópico e o teste de Wilcoxon pareado ($p = 0{,}717$) também não identificaram diferenças significativas. Assim, no contexto analisado, não há evidências empíricas de que o ChatGPT tenha alterado de forma mensurável a compreensão avaliada pelo instrumento utilizado.

Esses achados convergem com estudos que também não identificaram impacto significativo do ChatGPT no desempenho de estudantes em disciplinas de programação e estruturas de dados \cite{kosar2024chatgpt}. Ao mesmo tempo, a interpretação dos resultados deve considerar limitações importantes. A amostra final compreendeu 20 participantes, após cinco desistências; o questionário avaliou principalmente compreensão conceitual por meio de questões objetivas; e a alta familiaridade prévia dos estudantes com IA generativa, especialmente em programação, pode ter atenuado diferenças entre as condições experimentais.

Do ponto de vista educacional, a ausência de diferença significativa não deve ser interpretada como prova de neutralidade do ChatGPT para a aprendizagem. A ferramenta pode influenciar autonomia, dependência e engajamento cognitivo de formas não capturadas pelo instrumento adotado \cite{kosmyna2025brain,delikoura2025superficial}. Ainda assim, os resultados reforçam a necessidade de cautela ao assumir que o uso do ChatGPT em atividades práticas, por si só, altera a compreensão dos conteúdos de Heap e Tabelas Hash no curto prazo.

Como desdobramentos futuros, sugere-se ampliar o tamanho amostral, empregar instrumentos que contemplem diferentes dimensões da aprendizagem, investigar efeitos de longo prazo e comparar estratégias de mediação docente no uso da ferramenta. Também seria relevante replicar o estudo em outros tópicos de Algoritmos e Estruturas de Dados, ampliando a compreensão sobre como o ChatGPT influencia a formação em Computação em contextos educacionais distintos.

